%% kapa: exam=1 round=31 year=2020 period=1 form=A subject=civil_law subject_ko=민법 seq=18 number=18
%% source: Q-Net 감정평가사 기출문제 (cst00309) artlSeq=5207219 file=extracted/1008/5207219/A형/A형/2020년도 제31회 감정평가사 1차 시험_1교시_A형.pdf
%% conversion: kordoc 4.24.51 -> tools/gen1.py (auto); review=proofread-sonnet; flags=none
\begin{question}{18}
소멸시효에 관한 설명으로 옳지 않은 것은? (다툼이 있으면 판례에 따름)
\begin{choices}
  \item 인도받은 부동산을 소유권이전등기를 하지 않고 제3자에게 처분․인도한 매수인의 등기청구권은 소멸시효에 걸리지 않는다.
  \item 채무불이행으로 인한 손해배상청구권의 소멸시효는 손해배상을 청구한 때부터 진행한다.
  \item 채권자가 보증인을 상대로 이행을 청구하는 소를 제기한 때에도 주채무의 소멸시효가 완성하면 보증인은 주채무가 시효로 소멸되었음을 주장할 수 있다.
  \item 재산권이전청구권과 동시이행관계에 있는 매매대금채권의 소멸시효는 지급기일부터 진행한다.
  \item 등기 없는 점유취득시효가 완성하였으나 등기하지 않은 토지점유자가 토지의 점유를 잃은 경우, 그로부터 10년이 지나면 등기청구권은 소멸한다.
\end{choices}
\end{question}
\answer{2}
